\section{Introduction}

Modern GPU kernels are written against instructions that prescribe how their operands are laid out:
which thread holds which matrix element in which register, how a tile sits in shared memory,
which offsets a tensor-memory load touches. CuTe~\cite{cutedocs,cecka2026}, the core library of
CUTLASS~3 and~4 and of the CuTe DSL, represents all of these by one kind of object, the
\emph{layout}: a pair of nested integer tuples $\lay{\text{shape}}{\text{stride}}$, read as the
function that expands an index in the mixed radix given by the shape and takes the inner product
of the digits with the strides. Layouts are combined by an algebra---\code{coalesce},
\code{composition}, \code{complement}, \code{logical\_divide}, \code{logical\_product} and
relatives---with which a kernel author tiles a matrix, hands each thread its fragment, or checks
that a data layout matches an instruction. The same abstraction underlies
FlashAttention's kernels and several compilers.

The algebra is partial, and its definedness is subtle. \code{composition(A,B)} should be a layout
whose function is $A\circ B$; such a layout usually does not exist, the implementation decides
by ``stride'' and ``shape'' divisibility checks, and it reduces a multi-mode $B$ to its modes by a
distributive rule. The existing mathematical accounts are of three kinds. The author's
specification~\cite{cecka2026} states post-conditions and derives \emph{sufficient} conditions,
without proofs. Shah's note~\cite{shah2024} proves sufficient admissibility conditions for
complement and composition, and points out that the distributive rule is unsound in general.
Carlisle, Shah, Stern and VanKoughnett~\cite{carlisle2026} give a thorough treatment of flat
layouts and a categorical framework that explains the operations for the class of
\emph{tractable} layouts. None of them says exactly \emph{when} a composite exists, or why the
conditions and formulas are what they are; and, as we report below, the implementations do not
agree with the functional specification in either direction.

\paragraph{The idea.} A shape is a mixed-radix numeral system, and digit expansion is not additive:
adding two numbers produces carries. We show that this is the only thing that happens. Every
layout function can be written
\[ f(x)=x\cdot d_0+\sum_i\floor{x/P_i}\cdot\jump_i,\qquad \jump_i=d_i-s_{i-1}d_{i-1}, \]
with $P_i$ the prefix products of the shape, and satisfies the \emph{carry identity}
\[ f(x+y)=f(x)+f(y)+\sum_i c_i(x,y)\cdot\jump_i , \]
where $c_i(x,y)\in\{0,1\}$ is the carry into digit $i$. The ``jumps'' $\jump_i$ measure how far the
strides are from those of a linear function. Everything else is bookkeeping of carries.

\paragraph{Contributions.}
\begin{itemize}[leftmargin=1.4em,itemsep=2pt]
\item \textbf{A one-identity foundation} (\cref{sec:carry}), valid for strides in an arbitrary
 abelian group: integers of either sign, coordinate vectors, bit vectors under XOR.
\item \textbf{Which functions are layouts} (\cref{sec:normalform}). Every function on
 $\{0,\dots,N-1\}$ vanishing at $0$ has a unique expansion in the functions $\floor{x/Q}$; it is a
 layout function if and only if the scales $Q\ge2$ that occur form a divisibility chain dividing
 $N$ (\cref{thm:spectrum}). This re-proves, for all abelian groups, the theorem of
 \cite{carlisle2026} that \code{coalesce} is a normal form, and gives an $O(N\log N)$
 recognition algorithm; for mappings given by their values this answers the question,
 left open in~\cite[\S3.7]{bhaskaracharya2025}, of recovering shape and strides from a mapping.
\item \textbf{Exact composability} (\cref{sec:composition}). A composite, if it exists, is
 forced: its strides are the values of the outer layout at the strides of the inner one. It exists
 if and only if a \emph{carry defect} $\sum_i C_i\jump_i$ vanishes (\cref{thm:defect,thm:exact}).
 For outer layouts with independent jumps---for instance row-major, column-major and padded
 matrices---this is an explicit inequality (\cref{thm:carryfree}). CuTe's two divisibility
 conditions describe the case in which every stride of the inner layout has a single non-zero
 digit that divides the size of its mode, or whose multiples stay inside that mode
 (\cref{prop:monomial}); we prove the single-mode rule sound, and identify the two ways in
 which rule and truth differ: \emph{interference} between modes, which makes the distributive
 rule unsound, and composites with sums of strides, which the rule cannot produce.
\item \textbf{Complement and division} (\cref{sec:complement}). Complements are the direct-sum
 decompositions of an interval, a classical subject (Krasner factorisations;
 de~Bruijn~\cite{debruijn1956}, Long~\cite{long1967}) that long predates GPUs; this explains the formula and shows that the admissibility
 condition of \cite{shah2024,carlisle2026} cannot be weakened by any device. For composition with a compact layout---in particular
 for \code{logical\_divide}---we give a criterion that is exact for an arbitrary outer layout
 (\cref{thm:division}).
\item \textbf{The relation to XOR-based layouts and to the categorical framework}
 (\cref{sec:beyond}): Triton's linear layouts~\cite{zhou2026} are the layouts with values in
 $\F_2^k$ on binary shapes, and their composition is total because their intermediate addition
 has no carries; the category of~\cite{carlisle2026}, up to the bookkeeping of unit modes, sits
 inside a category of ``boxes'' in which
 composition is a lifting problem.
\item \textbf{Swizzles} (\cref{sec:swizzle}). Offsets carry two group laws, integer addition and
 XOR, which differ by the carries of the binary radix. A swizzle is a layout for XOR, and a
 swizzled layout is XOR-linear up to the parity of binary carries (\cref{thm:swzcomp}). We
 characterise when, for offset zero, it is an $\F_2$-valued layout, the kind of object used by
 PyCuTe and Triton
 (\cref{cor:xorlinear}); give a closed-form criterion, for bit-disjoint layouts, for a slice of
 it to decay to an ordinary layout with possibly negative strides (\cref{thm:decay}); and give
 conditions under which CuTe's operations on swizzled layouts are valid (\cref{sec:swzops}).
\item \textbf{Validation} (\cref{sec:mechanization}). The core of the theory---the floor-sum form,
 the carry identity, the characterisation of layouts as carry-free additive maps and by chain
 spectra, the composition formula with its defect, and the carry-free criterion---is formalised in Lean~4 with Mathlib
 (statements marked \mech); the main criteria and the examples are in addition checked by
 exhaustive and randomised programs; and differential tests compare the theory with three implementations (C++
 CuTe, CUTLASS's \code{pycute}, and the reference implementation PyCuTe), including their
 swizzle operations.
\end{itemize}

\paragraph{What is and is not new.} The normal-form theorem for coalesce, the existence and
uniqueness of complements under the divisibility condition, and a composition algorithm for
tractable layouts are due to~\cite{shah2024,carlisle2026}; sufficient conditions for composition
are in~\cite{shah2024,cecka2026}. New here are the carry identity and the slope--jump
coordinates as an organising principle, the spectral characterisation of layout functions, the
necessary and sufficient conditions for composition and division, the link to the classical
theory of direct-sum decompositions of intervals, proofs of these criteria for strides in an
arbitrary abelian group, and the mechanisation of the core.
Swizzles as layouts with XOR strides are due to~\cite{cecka2026}; their relation to CuTe's composed
layouts, the decay criterion and the conditions of \cref{sec:swzops} are new. We restrict
attention to the \emph{mathematics} of layouts; static versus dynamic integers, tensors, and code
generation are outside our scope.
